\documentclass[11pt,aspectratio=169]{beamer}

\usepackage[utf8]{inputenc}
\usepackage[T1]{fontenc}
\usepackage{lmodern}
\usepackage[spanish,es-noindentfirst]{babel}
\addto\captionsspanish{\renewcommand{\tablename}{Tabla}}
\usepackage{graphicx}
\usepackage{booktabs}
\usepackage{array}
\usepackage{xcolor}
\usepackage{amsmath}
\usepackage{tikz}

\definecolor{navy}{HTML}{1F4E5F}
\definecolor{teal}{HTML}{2E6B7E}
\definecolor{terra}{HTML}{C1440E}
\definecolor{mutedgray}{HTML}{5C7278}
\definecolor{ice}{HTML}{DCE9EC}

\usetheme{default}
\setbeamertemplate{navigation symbols}{}
\setbeamercolor{structure}{fg=navy}
\setbeamercolor{title}{fg=white,bg=navy}
\setbeamercolor{frametitle}{fg=white,bg=navy}
\setbeamerfont{frametitle}{series=\bfseries}
\setbeamercolor{block title}{fg=white,bg=teal}
\setbeamercolor{block body}{fg=black,bg=ice}
\setbeamercolor{itemize item}{fg=terra}
\setbeamercolor{itemize subitem}{fg=teal}
\setbeamertemplate{itemize item}[circle]
\setbeamertemplate{itemize subitem}[circle]
\setbeamercolor{footline}{fg=mutedgray}
\setbeamerfont{footline}{size=\tiny}
\setbeamertemplate{footline}{%
  \hfill\raisebox{2pt}{\color{mutedgray}\scriptsize\insertframenumber\,/\,\inserttotalframenumber}\hspace{4mm}\vspace{2mm}
}

\newcommand{\code}[1]{\texttt{\footnotesize #1}}

\title{\textbf{Más allá de la ideología}}
\subtitle{Tipologías empíricas de los partidos y movimientos políticos colombianos (1848--2023)}
\author{Andrez Felipe Guerrero Torres}
\institute{Aprendizaje de Máquinas y Políticas Públicas \\ Bogotá Summer School in Economics 2026 --- Pontificia Universidad Javeriana}
\date{}

\begin{document}

% ============================================================
\begin{frame}[plain]
\titlepage
\end{frame}

% ------------------------------------------------------------
\begin{frame}{Pregunta y motivación}
\begin{center}
\large\color{teal}\textquestiondown{}Qué grupos o patrones de partidos y movimientos políticos colombianos pueden identificarse mediante aprendizaje no supervisado (1848--2023)?
\end{center}
\vspace{0.4cm}
\begin{itemize}
  \item Cerca del \textbf{40\%} de los partidos no tiene ideología clasificable --- y esa ausencia no se concentra en los marginales.
  \item La literatura (Panebianco 1988; Mainwaring \& Scully 1995; Jones \& Mainwaring 2003) ofrece dimensiones más robustas: \textbf{institucionalización} (sobrevivir en el tiempo) y \textbf{nacionalización} (competir en múltiples niveles).
  \item Pizarro Leongómez (2006): el sistema colombiano, un ``gigante con pies de barro''.
\end{itemize}
\end{frame}

% ============================================================
% 01 -- DATOS
% ------------------------------------------------------------
\begin{frame}{01 \textbullet\ Los datos}
\begin{block}{Fuentes}
Cabra-Ruíz et al.\ (2023): 5.143 partidos/movimientos/coaliciones, 27 variables (1848--2023) \\
Torres et al.\ (2023): resultados electorales históricos (1958--2023), CEDE Uniandes
\end{block}
\begin{itemize}
  \item Códigos centinela (98/99) tratados como categorías, no imputados
  \item Votos transformados con $\log(1+\text{votos})$; coaliciones excluidas $\rightarrow$ 2.155 partidos no-coalición
  \item Dummies + estandarización para pesos iguales
  \item La participación cae con el nivel: 55,3\% en alcaldías vs.\ 1,2\% en presidencia
\end{itemize}
\end{frame}

% ============================================================
% 02 -- METODOLOGIA
% ------------------------------------------------------------
\begin{frame}{02 \textbullet\ K-Means y sus hiperparámetros}
\begin{columns}
\column{0.55\textwidth}
\begin{itemize}
  \item \textbf{Escala:} 5.143 obs.\ --- jerárquico poco práctico, K-Means escala linealmente
  \item \textbf{Datos dummy + euclidiana:} simple e interpretable
  \item \textbf{Validez verificable:} codo / silhouette / Calinski-Harabasz, a diferencia de DBSCAN o mezclas gaussianas
\end{itemize}
\column{0.45\textwidth}
\begin{center}
\begin{tabular}{ll}
$k$ (nivel 1) & 4 \\
$k$ (nivel 2) & 3 \\
\code{nstart} & 100 \\
\code{iter.max} & 1000 \\
Semilla & 20260706 \\
\end{tabular}
\end{center}
\end{columns}
\vspace{0.3cm}
\begin{block}{Importante}
Los nombres de los clusters son una interpretación teórica posterior, no una medición directa del algoritmo. PCA solo para visualización.
\end{block}
\end{frame}

% ------------------------------------------------------------
\begin{frame}{Selección de $k$ y modelo final}
\begin{columns}
\column{0.5\textwidth}
\includegraphics[width=\textwidth]{figuras/fig2_pca_clusters.png}
\column{0.5\textwidth}
\begin{itemize}
  \item Calinski-Harabasz favorece $k=2$; silhouette crece casi monótonamente
  \item Se eligió $k=4$: silhouette comparable a $k=2$--3 y partición interpretable
  \item Explica 25,2\% de varianza; silhouette promedio 0,26
  \item Cluster 1 (pequeño, clave): silhouette $-0,04$ --- interpretar con cautela
\end{itemize}
\end{columns}
\end{frame}

% ------------------------------------------------------------
\begin{frame}{Validación estadística}
\begin{columns}
\column{0.5\textwidth}
\textbf{Univariada} ($\chi^2$ + V de Cramér, las 12 variables $p<0,001$):
\begin{center}
\small
\begin{tabular}{lr}
\code{part\_presidencia} & 0,821 \\
\code{ideologia} & 0,615 \\
\code{gradonac} & 0,542 \\
\end{tabular}
\end{center}
\column{0.5\textwidth}
\textbf{Multivariada} (árbol + Random Forest, exactitud 99,1\%/99,6\%):
\end{columns}
\vspace{0.3cm}
\begin{block}{Hallazgo revelador}
\code{part\_presidencia} es la variable más fuerte aislada (V=0,821) pero cae al 5\textsuperscript{o} lugar en Random Forest: su señal es redundante con \code{ideologia} y \code{grupo\_representativo} --- algo que la prueba univariada no detecta.
\end{block}
\end{frame}

% ------------------------------------------------------------
\begin{frame}{Segundo nivel: votos reales y robustez de $k$}
\begin{columns}
\column{0.45\textwidth}
\includegraphics[width=\textwidth]{figuras/fig4_nivel2_pca.png}
\column{0.55\textwidth}
\begin{itemize}
  \item 2.155 partidos no-coalición, votos reales (1958--2023)
  \item Institucionalizado (4,8\%): Liberal, Conservador, Cambio Radical, Polo, Verde
  \item \textbf{Robustez:} $k=4$ replica casi intactos los clusters institucionalizado/nicho; Rand ajustado $=0,507$ entre $k=3$ y $k=4$
  \item Longevidad (externa, no usada en el clustering): 13,3 años institucionalizado vs.\ 0,5--0,8 en los demás
\end{itemize}
\end{columns}
\end{frame}

% ============================================================
% 03 -- OTRAS METODOLOGIAS
% ------------------------------------------------------------
\begin{frame}{03 \textbullet\ Otras metodologías (chequeos de robustez)}
\begin{block}{Kernel PCA + Spectral Clustering}
Núcleo RBF, muestra de 2.000 partidos: acuerdo moderado con K-Means (Rand ajustado $=0,57$); no aísla el 1,8\% institucionalizado como grupo propio --- perfil sustantivo pero sensible al método.
\end{block}
\begin{block}{Naive Bayes sobre el nombre del partido}
Exactitud 84,1\% vs.\ 85,7\% de la clase mayoritaria (baseline) --- resultado negativo: el nombre no predice de forma confiable la institucionalización.
\end{block}
\end{frame}

% ============================================================
% 04 -- CONCLUSIONES
% ------------------------------------------------------------
\begin{frame}[plain]
\begin{tikzpicture}[remember picture,overlay]
\fill[navy] (current page.south west) rectangle (current page.north east);
\end{tikzpicture}
\centering
\vspace{0.9cm}
{\color{white}\large 04 \textbullet\ Hallazgo central\par}
\vspace{0.4cm}
{\color{terra}\fontsize{46}{54}\selectfont\bfseries 1,8\%\par}
\vspace{0.2cm}
{\color{white}\normalsize 93 de 5.143 partidos institucionalizados y nacionalizados.\\ El 98,2\% restante opera de forma local o marginal.\\ El segundo nivel corrobora: 4,8\% institucionalizado.\par}
\vspace{0.5cm}
{\color{ice}\footnotesize Las reformas de 2003 y 2011 no lo frenaron: la probabilidad de institucionalizarse cayó de 14,4\% a 1,2\%.\par}
\end{frame}

% ------------------------------------------------------------
\begin{frame}{Conclusiones de política pública}
\begin{itemize}
  \item \textbf{Personería jurídica:} exigir desempeño multinivel \textbf{sostenido}, no solo una elección puntual
  \item \textbf{Financiación estatal:} indicadores de presencia multinivel, no solo volumen de votos
  \item \textbf{Partidos programáticos:} la débil asociación institucionalización-ideología exige instrumentos específicos
  \item \textbf{Monitoreo:} la tipología, aplicada periódicamente por Registraduría/CNE, para evaluar si las reformas consolidan o fragmentan el sistema
\end{itemize}
\vspace{0.3cm}
\begin{block}{Advertencia}
Instrumento analítico de apoyo, \textbf{no} un criterio automático para cancelar personerías jurídicas.
\end{block}
\vspace{0.3cm}
\centering \textexclamdown{}Gracias! \hspace{0.3cm} \url{https://github.com/andrezconz/curso-machine-learning-2026-tf}
\end{frame}

\end{document}
